\section{Stage 243: Relaxed constraint branch declaration}
\label{sec:app-part08:relaxed-constraint-branch-declaration}
\label{stage:243}
\stagefield{Verification}{SymPy audit: \StageFile{scripts/moving_throat_pde_stage243_relaxed_constraint_branch_declaration_and_short_range_open_system_compiler_sympy_audit.py}.  Mathematica audit: \StageFile{mathematica/moving_throat_pde_stage243_relaxed_constraint_branch_declaration_and_short_range_open_system_compiler_mathematica_audit.wl}.}


Stage 243 declares the relaxed branch without yet assembling a barrier.  Its purpose is to make the Session-I relaxations legal inside the derivation tree while preserving the same-charge short-range verdict.  The inherited base object is the Stage 242 front-end packet
\begin{equation}
\label{eq:app-part08-p225-packet}
\mathcal P_{242}
=
\Bigl(
\epsilon,
\varrho_{\rm phys},
\sigma_{\rm phys},
\text{ranking region},
R_{\rm tr},
R_{\rm target},
\epsilon_\eta,
 d\ln R_{\rm tr},
 d\ln R_{\rm target},
 d\ln\epsilon_\eta,
 N_Q-1
\Bigr),
\end{equation}
with selected-support packet
\begin{equation}
\label{eq:app-part08-s225-packet}
\mathcal S_{242}=(\zeta,M_{\rm mix},S(\zeta;\epsilon),M_{\rm tr}).
\end{equation}
The strict slice is
\begin{equation}
\label{eq:app-part08-std-branch-slice}
\mathfrak B_{\rm std}
=
\left\{
J^w=0,
\;U=0,
\;V=0,
\;\varsigma(z)\equiv 1
\right\}
\subset(\mathcal P_{242},\mathcal S_{242}).
\end{equation}
This definition is useful because it names exactly what is being relaxed: transverse-current suppression, rigid-mouth orbit lock, and positive-source Family-1 closure.

\subsection{Open-system leakage/work lane}
\label{subsec:app-part08-open-lane}

The open-system lift begins with the exact projected-continuity leakage identity
\begin{equation}
\label{eq:app-part08-stage243-leakage-identity}
S_{\rm leak}
=
-\bigl[W(w)j^w\bigr]_{-\infty}^{+\infty}
+
\int_{-\infty}^{+\infty}W'(w)j^w(w)\,dw.
\end{equation}
For the declaration-level Gaussian closure,
\begin{equation}
\label{eq:app-part08-stage243-gaussian-profile}
W(w)=\frac{e^{-w^2}}{\sqrt\pi},
\qquad
j^w(w)=\ell_w j_0\,w e^{-w^2},
\qquad
E_w(w)=E_0we^{-w^2},
\end{equation}
so the boundary term vanishes and
\begin{equation}
\label{eq:app-part08-stage243-leak-work}
S_{\rm leak}=-\frac{\sqrt2}{4}\ell_wj_0,
\qquad
\mathcal W_w=\frac{\sqrt{2\pi}}{8}\ell_wj_0E_0.
\end{equation}
The declaration packet for this lane is
\begin{equation}
\label{eq:app-part08-stage243-lw}
L_w=(S_{\rm leak},\mathcal W_w).
\end{equation}
Both components vanish on \(\ell_w=0\).

\subsection{Non-rigid mouth/dressing lane}
\label{subsec:app-part08-uv-lane}

The non-rigid lift uses the minimal free energy
\begin{equation}
\label{eq:app-part08-stage243-fuv}
\mathcal F_{UV}(U,V)
=
\frac12k_UU^2+\frac12k_VV^2-\chi_\lambda UV-f_UU.
\end{equation}
Stationarity gives
\begin{equation}
\label{eq:app-part08-stage243-uv-solution}
U=\frac{k_Vf_U}{k_Uk_V-\chi_\lambda^2},
\qquad
V=\frac{\chi_\lambda f_U}{k_Uk_V-\chi_\lambda^2},
\qquad
\frac{V}{U}=\frac{\chi_\lambda}{k_V}.
\end{equation}
The branch is admissible when
\begin{equation}
\label{eq:app-part08-stage243-uv-stability}
k_Uk_V-\chi_\lambda^2>0,
\end{equation}
and the drain is
\begin{equation}
\label{eq:app-part08-stage243-duv}
\mathcal D_{UV}=\chi_\lambda UV
=
\frac{\chi_\lambda^2k_Vf_U^2}{(k_Uk_V-\chi_\lambda^2)^2}\ge0.
\end{equation}
This lane is the finite relaxed continuation of the earlier weak-axisymmetric transfer-shape scalar: at infinitesimal order, \(\Xi_1=\delta U\).

\subsection{Compensated sign-changing source lane}
\label{subsec:app-part08-compensated-source-lane}

The compensated source lift is
\begin{equation}
\label{eq:app-part08-stage243-varsigma}
\varsigma(z)=1+a\cos(\pi z)+b\cos(2\pi z),
\qquad
z\in[0,1],
\end{equation}
with exact unit mean.  Setting \(y=\cos(\pi z)\),
\begin{equation}
\label{eq:app-part08-stage243-varsigma-quadratic}
\varsigma(y)=1-b+ay+2by^2,
\qquad
-1\le y\le 1.
\end{equation}
The interior stationary point and value are
\begin{equation}
\label{eq:app-part08-stage243-varsigma-min}
y_*=-\frac{a}{4b},
\qquad
\varsigma(y_*)=1-b-\frac{a^2}{8b},
\end{equation}
with the boundary candidates \(1+a+b\) and \(1-a+b\).  Thus the first source relaxation is not a vague sign-changing profile; it is an exact mean-preserving quadratic problem.

The full declaration packet is
\begin{equation}
\label{eq:app-part08-stage243-relaxed-branch}
\mathfrak B_{243}^{\rm relax}
=
\left\{(\mathcal P_{242},\mathcal S_{242});\;L_w,\;L_{UV},\;L_\varsigma\right\},
\end{equation}
where
\begin{equation}
\label{eq:app-part08-stage243-lanes}
L_w=(S_{\rm leak},\mathcal W_w),
\quad
L_{UV}=(U,V,\mathcal D_{UV}),
\quad
L_\varsigma=(\varsigma,\varsigma_{\min},\mathrm{signchg}).
\end{equation}

\subsection{Short-range theorem for the lifted branch}
\label{subsec:app-part08-short-range-theorem}

The most important Stage 243 output is the short-range firewall.  With primitive profiles
\begin{equation}
\label{eq:app-part08-stage243-source-profiles}
\mathcal S_Q(r)=r^{-3},
\qquad
\mathcal S_Y(r)=\frac{e^{-2\kappa r}}{r},
\end{equation}
the static one-port correction is
\begin{equation}
\label{eq:app-part08-stage243-static-correction}
\delta V_{\rm stat}(r)
=
-\frac12\left[
\frac{\mathcal C_6}{r^6}
+2\mathcal C_4\frac{e^{-2\kappa r}}{r^4}
+\mathcal C_2\frac{e^{-4\kappa r}}{r^2}
\right],
\end{equation}
and the linear monochromatic dynamic bundle has the same spatial span:
\begin{equation}
\label{eq:app-part08-stage243-dynamic-correction}
\Re\mathfrak V_{\rm dyn}(r,\omega)
=
-\frac12\left[
\frac{\mathcal C_6(\omega)}{r^6}
+2\mathcal C_4(\omega)\frac{e^{-2\kappa r}}{r^4}
+\mathcal C_2(\omega)\frac{e^{-4\kappa r}}{r^2}
\right].
\end{equation}
Thus the relaxed branch is a short-range/open-system bypass branch.  It can reshape near-contact energetics, but it cannot be cited as a new long-range same-charge force.
